\documentclass{article}%
\usepackage[T1]{fontenc}%
\usepackage[utf8]{inputenc}%
\usepackage{lmodern}%
\usepackage{textcomp}%
\usepackage{lastpage}%
\usepackage{geometry}%
\geometry{margin=2.2cm}%
%
\usepackage{amsmath}%
\usepackage{amssymb}%
\usepackage{booktabs}%
\usepackage{longtable}%
\usepackage{array}%
\usepackage{hyperref}%
\hypersetup{hidelinks}%
\title{IME 2027 {-} Física (claude\_sonnet\_5\_5) x Poliedro\\[4pt]\large Comparação com o gabarito Poliedro\thanks{Gerado com gpt-resolve: \url{https://github.com/lgabs/gpt-resolve}}}%
\date{09/11/2026}%
%
\begin{document}%
\normalsize%
\maketitle%
\section*{Resumo}\textbf{Nota média: 10,0/10} (30 de 30 itens com a resolução correta; cada item vale o mesmo dentro da questão).\par\medskip%
\begin{longtable}{@{}c c p{0.78\textwidth}@{}}\toprule Questão & Nota & Resumo \\ \midrule
1 & 10,0 & A resolução está correta nos três itens e bate com o gabarito Poliedro: v = 250 m/s, T = 4000 K e W = 864 kJ (Q = 2880 kJ). Refiz as contas em python3 e conferi a transcrição com o enunciado original, que está fiel. A resolução ainda verifica que o foguete está subindo na vertical do alvo e confere o trabalho por um segundo caminho (W = m$\cdot$(R/M)$\cdot$$\Delta$T). Não encontrei erro em nenhum dos dois. \\
2 & 10,0 & A resolução e o gabarito Poliedro coincidem nos três itens. a) A = 200 kN para cima e B = 100 kN para baixo, ambos verticais. b) CD 75 C, CF 125 T, DF 200 C, AD 200 C, AE 0, DE 75 C, BE 100 T e EF 125 T (kN). c) AE é barra nula neste carregamento, mas funciona como contraventamento contra o mecanismo do paralelogramo ADEB. Conferi a figura no PDF original e resolvi o sistema completo da treliça com frações exatas: os valores batem e a estrutura é isostática (posto 12). Sem AE, vira um mecanismo (posto 11 para 12 equações). A única ressalva é a redação confusa de uma linha de conferência na resolução, que não afeta a nota. \\
3 & 10,0 & Os três itens da resolução coincidem com o gabarito Poliedro, e as contas conferidas em python estão certas. a) R = 300 N para cima (T = 300 N). b) $\Delta$T = 100 $^\circ$C. c) Q = 5$\cdot$10\^{}6 J e rompimento em 37,5\% do dia. O enunciado lido no PDF original confere com a transcrição. Ficam só hipóteses simplificadoras de baixo impacto, comuns às duas resoluções: área de absorção fixa e nenhuma perda de calor. \\
4 & 10,0 & A resolução está correta nos três itens e coincide com o gabarito Poliedro. a) V\_AB = 900 V, i1 = 50 A e i2 = 0. b) A razão é 2,5, acima da assíntota de cerca de 1,1, e o fusível rompe em cerca de 1,5 s; a medição em pixels da Figura 2 confirma 1,50 s. c) Com o curto, I = 500 A e a razão é 25; o patamar da curva dá 100 ms, mais que os 50 ms suportáveis, então o fusível não protege o gerador. Não há divergência. Os únicos pontos a observar são a leitura gráfica do item (b) e a extrapolação do patamar até a razão 25 no item (c). \\
5 & 10,0 & A resolução está correta e bate com o gabarito Poliedro nos três itens: a) T = 450 K; b) Q = 90 kW; c) as alegações de B não são realistas. As 350 residências exigem 116,7 kW, mas com a água aquecida (33,3 kW) e os 45 kW rejeitados a 450 K, B entrega no máximo 71,7 kW, cerca de 215 residências. A resolução também acrescenta o limite de 85 kW dado pela 2ª lei para o caso sem água, e essa conta está certa. Conferi todas as contas em python. Nada no gabarito nem na resolução muda a nota. \\
6 & 10,0 & A resolução e o gabarito Poliedro concordam nos quatro itens: v = 68 m/s, R = 5$\surd$3 $\approx$ 8,66 m, T $\approx$ 1,07$\times$10$^3$ N e V\_sub $\approx$ 0,055 m$^3$ (E $\approx$ 551 N). Refiz as contas em Python e as quatro conferem. A física está correta: Doppler com a componente radial da velocidade no instante da emissão, sen $\theta$ = R/L para o fio, componente horizontal da tração como força centrípeta e equilíbrio vertical E = mg + T$\cdot$cos $\theta$. As diferenças são só de notação e de nome dos ângulos. Não encontrei erro na resolução nem no gabarito. \\
7 & 10,0 & A resolução está correta nos três itens e bate com o gabarito Poliedro: a) d = 0,45 m; b) kx $\approx$ 6,44 N e razão de módulo 29/18 $\approx$ 1,6; c) v = $\surd$2 m/s ao atingir C. Refiz todas as contas em python3 e confirmei os valores. A única diferença é o sinal no item b. O gabarito dá -1,6, que é a razão literal, já que $\Delta$U\_g < 0. A resolução dá +1,6 usando |$\Delta$E\_pg| de forma explícita, o que é coerente mas não traz o sinal pedido pela palavra \textbackslash{}"variação\textbackslash{}". Se o sinal for exigido, o item b pode perder pontos. \\
8 & 10,0 & A resolução e o gabarito Poliedro concordam nos quatro itens. Os dois usam o mesmo modelo: d = 2,5 µm e sen $\theta$\_m = m/4, e o máximo é visível enquanto D$\cdot$tg$\theta$\_m $\leq$ 0,6 m. a) O anteparo está se afastando da rede. b) D = 0,2$\surd$7 $\approx$ 0,53 m. c) v $\approx$ 0,254 m/s. d) $\Delta$s = 0,6$\surd$15 - 0,2$\surd$7 $\approx$ 1,79 m. Refiz as contas em python3 e elas conferem. As únicas diferenças vêm de arredondamento ou do uso da tabela de raízes. \\
9 & 10,0 & A resolução e o gabarito Poliedro concordam nos dois itens. No item (a), a fórmula é a mesma: u = c($\surd$3h - L)/($\surd$3L - h). No item (b), h < L < $\surd$3h é a mesma condição que L/$\surd$3 < h < L, só escrita de outro jeito. Refiz a conta e conferi numericamente em Python que o raio refletido sai na direção (0, -1) e que o impacto é o primeiro contato. Não há erro em nenhum dos dois. Conferi a transcrição com o PDF e ela está fiel ao original. \\
10 & 10,0 & A resolução e o gabarito do Poliedro estão certos e coincidem nos dois itens. Em (a), a resposta é somente a situação 2, com a mesma justificativa para descartar as situações 1, 3 e 4. Em (b), p $\approx$ ($\pi$d/2)$\surd$($\pi$mv/($\mu$0 i q)), que corresponde a T = $\pi$d$\surd$($\pi$m/($\mu$0 i q v)) e p = vT/2. Refiz a conta, que confere com o enunciado original (p. 10). Uma simulação numérica com o campo exato dos dois fios confirma a fórmula com diferença de cerca de 0,06\%. Não encontrei erro em nenhum dos dois. \\
\bottomrule\end{longtable}%
\section*{Principais pontos de dúvida}%
\begin{itemize}%
\item%
\textbf{Q7} \textit{(solution)} Item b: o enunciado pede a razão entre U\_el e a "variação da energia potencial gravitacional de P2". Essa variação é negativa, porque P2 desce. O gabarito responde $\cong$ -1,6. A resolução responde +1,6 e deixa explícito que usa |$\Delta$E\_pg|. As duas contas são coerentes entre si, mas um corretor rigoroso pode descontar pela falta do sinal. Vale levar ao professor para saber se o sinal é exigido.%
\end{itemize}%
\section*{Questão 1}%
\href{https://poliedroresolve.sistemapoliedro.com.br/vestibulares/ime/2027/ime-2-fase-dia-2-27-10-2026-2-fase-dia-2/dissertativa-1-fisica-2-fase-dia-2-ime-2027}{Resolução do gabarito}\par\smallskip%
\begin{longtable}{@{}p{0.06\textwidth} p{0.33\textwidth} p{0.33\textwidth} p{0.1\textwidth} p{0.08\textwidth}@{}}\toprule Item & Resolução & Gabarito & Veredito & Certo \\ \midrule
a & v = 250 m/s & v0 = 250 m/s & igual & ambos \\
b & T = 4000 K & T\_falha = 4000 K & igual & ambos \\
c & W = 864 kJ (realizado pelo gás; Q = 2880 kJ, $\Delta$U = 2016 kJ) & W = 864 kJ (Q = 2880 kJ, $\Delta$U = 2016 kJ) & igual & ambos \\
\bottomrule\end{longtable}%
\begin{itemize}%
\item%
\textbf{a)} A equação da trajetória com $\theta$=45$^\circ$, x=2000 m e y=1360 m dá v$^2$ = g$\cdot$x$^2$/(x-y) = 10$\cdot$4$\cdot$10$^{6}$/640 = 62500, logo v = 250 m/s. A resolução também confirma que o foguete ainda sobe: vy $\approx$ 63,6 m/s > 0 e o ápice fica em x = 3125 m.%
\item%
\textbf{b)} Como T $\propto$ v$^2$, T = 2560$\cdot$(250/200)$^2$ = 2560$\cdot$1,5625 = 4000 K. Conta conferida.%
\item%
\textbf{c)} Pela 1ª Lei, com Q = m$\cdot$cp$\cdot$$\Delta$T = 2$\cdot$1$\cdot$1440 = 2880 kJ, W = Q - $\Delta$U = 2880 - 2016 = 864 kJ. A resolução confere por outro caminho: cv = 0,70, R/M = 0,30 kJ/(kg$\cdot$K) e m(R/M)$\Delta$T = 864 kJ. Também indica o sentido (trabalho feito pelo gás, que se expande).%
\end{itemize}%
\paragraph{Pontos de dúvida}%
\begin{itemize}%
\item%
\textit{(statement)} O enunciado trata o foguete como projétil lançado com velocidade inicial dada, sem fase propulsada, e admite T $\propto$ v$^2$ como hipótese. São idealizações do próprio enunciado, e tanto a resolução quanto o gabarito as seguem. Isso não afeta a nota.%
\item%
\textit{(statement)} No item c, "trabalho realizado durante o aquecimento" não diz se o trabalho é feito pelo gás ou sobre ele. A resolução declara a convenção (pelo gás, W > 0 na expansão isobárica), e ela coincide com a do gabarito ($\Delta$U = Q - W).%
\end{itemize}%
\section*{Questão 2}%
\href{https://poliedroresolve.sistemapoliedro.com.br/vestibulares/ime/2027/ime-2-fase-dia-2-27-10-2026-2-fase-dia-2/dissertativa-2-fisica-2-fase-dia-2-ime-2027}{Resolução do gabarito}\par\smallskip%
\begin{longtable}{@{}p{0.06\textwidth} p{0.33\textwidth} p{0.33\textwidth} p{0.1\textwidth} p{0.08\textwidth}@{}}\toprule Item & Resolução & Gabarito & Veredito & Certo \\ \midrule
a & A: 200 kN, vertical, para cima (A\_x = 0); B: 100 kN, vertical, para baixo (B\_x = 0) & A = 200 kN vertical para cima; B = 100 kN vertical para baixo (A\_x = B\_x = 0) & igual & ambos \\
b & CD = 75 (C); CF = 125 (T); DF = 200 (C); AD = 200 (C); AE = 0; DE = 75 (C); BE = 100 (T); EF = 125 (T), em kN & N\_CD = 75 kN (C); N\_CF = 125 kN (T); N\_DF = 200 kN (C); N\_AD = 200 kN (C); N\_AE = 0; N\_DE = 75 kN (C); N\_BE = 100 kN (T); N\_EF = 125 kN (T) & igual & ambos \\
c & AE tem esforço nulo para P vertical, mas é contraventamento: triangula o paralelogramo articulado ADEB, impede o mecanismo (oscilação lateral, equilíbrio instável), torna a treliça isostática e absorve cargas horizontais. & N\_AE = 0 só para este carregamento; AE é necessária para outras distribuições de carga e garante a rigidez (sem ela ADEB vira paralelogramo articulado e o guindaste balança lateralmente) & igual & ambos \\
\bottomrule\end{longtable}%
\begin{itemize}%
\item%
\textbf{a)} Refiz com o sistema completo do método dos nós (12 equações, 12 incógnitas, frações exatas): A\_x = 0, A\_y = +200 kN, B\_x = 0, B\_y = -100 kN. Também confere pelo momento em B: P$\cdot$6 = A\_y$\cdot$3 $\Rightarrow$ A\_y = 200 kN. B\_x = 0 porque a única barra que chega em B (BE) é vertical.%
\item%
\textbf{b)} As oito forças e suas naturezas batem com a solução exata do sistema: CD -75, DE -75, CF +125, DF -200, EF +125, AD -200, AE 0, BE +100 (positivo = tração). Conferi também as equações nó a nó da resolução (C, F, D, E), e os sinais dos versores estão corretos.%
\item%
\textbf{c)} A resolução diz o mesmo que o gabarito e acrescenta detalhes corretos. Sem AE, a matriz de equilíbrio tem 11 incógnitas e posto 11 para 12 equações, ou seja, um mecanismo de 1 grau de liberdade (7 + 4 = 11 < 12). A afirmação de que esse equilíbrio seria instável também está certa: num deslocamento lateral $\delta$, AD (compressão de 200) desestabiliza com 200$\delta$/4, e BE (tração de 100) restaura com 100$\delta$/4. O saldo de 25$\delta$ desestabiliza, como num pêndulo invertido.%
\end{itemize}%
\paragraph{Pontos de dúvida}%
\begin{itemize}%
\item%
\textit{(solution)} A linha 'Conferência dos nós' do item b está truncada e confusa ('em F, -75-75+\ldots{} fecha com +200 vertical de DF (-100-100+200=0)'). Mistura componentes horizontais e verticais de CF/EF, embora os números finais estejam certos. É só redação, não muda resultado.%
\item%
\textit{(solution)} No item c, a frase 'o nó E não tem componente horizontal a equilibrar' é imprecisa. Na verdade as componentes horizontais de DE (75 kN) e de EF (3/5$\cdot$125 = 75 kN) se anulam em E, e por isso AE = 0. A conclusão está certa.%
\item%
\textit{(statement)} A reação em B para baixo pressupõe que o apoio articulado fixo resista a arrancamento (esteja ancorado). É hipótese implícita no enunciado, e tanto a resolução quanto o gabarito a adotam.%
\end{itemize}%
\section*{Questão 3}%
\href{https://poliedroresolve.sistemapoliedro.com.br/vestibulares/ime/2027/ime-2-fase-dia-2-27-10-2026-2-fase-dia-2/dissertativa-3-fisica-2-fase-dia-2-ime-2027}{Resolução do gabarito}\par\smallskip%
\begin{longtable}{@{}p{0.06\textwidth} p{0.33\textwidth} p{0.33\textwidth} p{0.1\textwidth} p{0.08\textwidth}@{}}\toprule Item & Resolução & Gabarito & Veredito & Certo \\ \midrule
a & R = 300 N, vertical, para cima (T0 = 300 N; componente horizontal nula) & T = 300 N; R\_A = 300 N vertical para cima & igual & ambos \\
b & $\Delta$T = 100 $^\circ$C & $\Delta$T = 100 $^\circ$C & igual & ambos \\
c & Rompe em 37,5\% do dia ($\tau$ = 10\^{}4 s = 12,5\% do dia após o nascer do Sol), com Q = 5$\cdot$10\^{}6 J & Q = 5$\cdot$10\^{}6 J; rompe às 37,5\% do dia (12,5\% após 25\%) & igual & ambos \\
\bottomrule\end{longtable}%
\begin{itemize}%
\item%
\textbf{a)} m = 2000$\cdot$(2$\cdot$0,25$\cdot$0,05) = 50 kg, P = 50$\cdot$12 = 600 N. Momentos na articulação: T$\cdot$L = P$\cdot$L/2 $\Rightarrow$ T = 300 N; R = P - T = 300 N para cima. Conferido numericamente.%
\item%
\textbf{b)} O ponto da argola fica fixo a L da articulação e o CM passa para L'/2 = L(1+$\alpha$$\Delta$T)/2. Então T = T0(1+$\alpha$$\Delta$T). Com T = 1,001$\cdot$T0, $\alpha$$\Delta$T = 10\^{}{-}3 e $\Delta$T = 100 $^\circ$C. As duas resoluções fazem o mesmo raciocínio.%
\item%
\textbf{c)} Q = mc$\Delta$T = 50$\cdot$10$^3$$\cdot$100 = 5$\cdot$10\^{}6 J, ou seja, Q/A = 10\^{}7 J/m$^2$ com A = 0,5 m$^2$. No trecho crescente, I = 0,2$\tau$ e a energia acumulada por área é 0,1$\tau$$^2$ = 10\^{}7, logo $\tau$ = 10\^{}4 s = 12,5\% do dia, antes do pico (o trecho crescente fornece 4$\cdot$10\^{}7 J/m$^2$). O horário é 25\% + 12,5\% = 37,5\%. O gabarito usa t em \% do dia e chega a 32000$\cdot$t$^2$ = 5$\cdot$10\^{}6, também t = 12,5\%. As duas formas são equivalentes.%
\end{itemize}%
\paragraph{Pontos de dúvida}%
\begin{itemize}%
\item%
\textit{(statement)} A dilatação superficial da face superior foi desprezada (A = 0,5 m$^2$ fixa). Com $\Delta$T = 100 $^\circ$C a área cresce até \textasciitilde{}0,2\%, o que altera o horário em cerca de 0,05\% do valor, como observa o Poliedro. O efeito é desprezível, mas é uma hipótese que nem a resolução nem o enunciado declaram com todas as letras.%
\item%
\textit{(statement)} Supõe{-}se que não há perdas de calor (radiação de um corpo escuro, condução pelos apoios) e que a temperatura da barra é uniforme. O enunciado só diz que a irradiância é 'totalmente absorvida'. É uma hipótese razoável e de praxe nessa questão.%
\end{itemize}%
\section*{Questão 4}%
\href{https://poliedroresolve.sistemapoliedro.com.br/vestibulares/ime/2027/ime-2-fase-dia-2-27-10-2026-2-fase-dia-2/dissertativa-4-fisica-2-fase-dia-2-ime-2027}{Resolução do gabarito}\par\smallskip%
\begin{longtable}{@{}p{0.06\textwidth} p{0.33\textwidth} p{0.33\textwidth} p{0.1\textwidth} p{0.08\textwidth}@{}}\toprule Item & Resolução & Gabarito & Veredito & Certo \\ \midrule
a & V\_AB = 900 V; i1 = 50 A (de A para B); i2 = 0; I = 50 A & U\_AB = 900 V; i1 = 50 A; i2 = 0 & igual & ambos \\
b & I/I\_FUS = 50/20 = 2,5 > 1,1 (assíntota): o fusível rompe em t $\approx$ 1,5 s (faixa de 1,4 a 1,6 s) & I/I\_fus = 2,5: o fusível rompe em t $\approx$ 1,5 s (lido na curva) & igual & ambos \\
c & I\_cc = 1000/2 = 500 A, razão 25; pelo patamar, t = 100 ms > 50 ms, logo o fusível NÃO protege o gerador & I\_curto = 500 A, I/I\_fus = 25; t = 100 ms > 50 ms, logo o fusível NÃO protege a fonte & igual & ambos \\
\bottomrule\end{longtable}%
\begin{itemize}%
\item%
\textbf{a)} Refiz a conta. No original, todas as baterias têm o polo + voltado para A. Assim, V = 1000 - 2I = 850 + i1 = 900 + 2i2, com I = i1 + i2. Por Millman, V = (500 + 850 + 450)/(0,5 + 1 + 0,5) = 900 V, o que dá i1 = 50 A, i2 = 0 e I = 50 A. A resolução e o gabarito coincidem, e a resolução ainda confere o balanço de potência (50 kW = 5 kW + 45 kW).%
\item%
\textbf{b)} Medi a Figura 2 em pixels na página renderizada a 300 dpi. O trecho inclinado vai de cerca de (1,08; 3,83 s) até cerca de (3,35; 0,1 s). Interpolando em 2,5, obtém{-}se t $\approx$ 1,50 s, e a reta passa por cerca de 1,5 s na linha de grade x = 2,5. O valor é o mesmo do gabarito e da resolução. A resolução justifica corretamente que 2,5 está acima da assíntota (cerca de 1,1), de modo que o rompimento ocorre.%
\item%
\textbf{c)} Com o curto entre A e B, V\_AB = 0. O fusível fica só no laço do gerador, então I = 1000/2 = 500 A e a razão é 25. Na curva, acima de cerca de 3,35 o tempo é constante em 10$^{-}$$^{1}$ s = 100 ms, que é maior que os 50 ms suportáveis. A conclusão de que o fusível não protege a fonte é a mesma nos dois. As informações extras da resolução também estão corretas: as correntes dos receptores não passam por F, e a energia dissipada no gerador é de 50 kJ, contra 25 kJ em 50 ms.%
\end{itemize}%
\paragraph{Pontos de dúvida}%
\begin{itemize}%
\item%
\textit{(statement)} A Figura 2 só vai até I/I\_FUS = 5. Para a razão 25 do item (c), tanto a resolução quanto o gabarito supõem que o patamar de 10$^{-}$$^{1}$ s continua indefinidamente. É a leitura natural, já que o fusível tem um tempo mínimo de atuação, mas é uma extrapolação. A resolução declara essa hipótese, e o gabarito a deixa implícita.%
\item%
\textit{(key)} O tempo de 1,5 s do item (b) vem da leitura de um gráfico. Pela minha medição em pixels, os extremos do trecho inclinado ficam em cerca de (1,08; 3,83 s) e (3,35; 0,10 s), o que dá t(2,5) $\approx$ 1,50 s. Dependendo de como se leem os extremos, valores entre cerca de 1,4 e 1,6 s devem ser aceitos.%
\item%
\textit{(statement)} O modelo usa fcem constantes mesmo 'imediatamente após o fechamento de S'. Num motor real parado, a fcem é nula na partida. O enunciado dá as fcem como valores fixos, e a resolução e o gabarito seguem esse modelo. Isso não altera a nota.%
\item%
\textit{(solution)} No item (c), a resolução afirma que os receptores também ficam em curto, com i1 = -850 A e i2 = -450 A. Isso decorre do modelo de fcem ideal, mas é irrelevante para o fusível e pouco físico se os receptores forem motores na partida. Não afeta a resposta.%
\end{itemize}%
\section*{Questão 5}%
\href{https://poliedroresolve.sistemapoliedro.com.br/vestibulares/ime/2027/ime-2-fase-dia-2-27-10-2026-2-fase-dia-2/dissertativa-5-fisica-2-fase-dia-2-ime-2027}{Resolução do gabarito}\par\smallskip%
\begin{longtable}{@{}p{0.06\textwidth} p{0.33\textwidth} p{0.33\textwidth} p{0.1\textwidth} p{0.08\textwidth}@{}}\toprule Item & Resolução & Gabarito & Veredito & Certo \\ \midrule
a & T = 450 K ($\eta$\_A = 0,40; $\eta$\_max = 0,50) & $\eta$1 = 40\%; T = 450 K & igual & ambos \\
b & Q = 90 kW (verificação: $\dot S$\_ger = 0,033 kW/K > 0) & Q rejeitado = 90 kW & igual & ambos \\
c & Não são realistas: 350 residências exigem 116,7 kW e a água 33,3 kW; somados aos 45 kW rejeitados a 450 K dão 195 kW > 150 kW (viola a 1ª lei). Com a água, $\dot W$\_B $\leq$ 71,7 kW, o que dá cerca de 215 residências. Sem água, a 2ª lei limita $\dot W$\_B a 85 kW (255 residências). O caso de 71,7 kW com a água é compatível com a 2ª lei (Q2 = 33,3 > 20 kW). & P = 116,7 kW; Q\_água = 33,3 kW; a máquina 2 fornece 71,7 kW < 116,7 kW $\Rightarrow$ NÃO supre as 350 residências (máx. 215); $\Delta$$\dot S$\_universo = +0,04 kW/K $\Rightarrow$ não viola a 2ª lei & igual & ambos \\
\bottomrule\end{longtable}%
\begin{itemize}%
\item%
\textbf{a)} $\dot W$ = 80$\cdot$3/4 = 60 kW; $\eta$ = 60/150 = 0,4 = 0,8(1 - T/900) $\Rightarrow$ T = 450 K. Conferido em python.%
\item%
\textbf{b)} Pela 1ª lei, 150 - 60 = 90 kW. A verificação de entropia da resolução (90/450 - 150/900 = 0,033 kW/K > 0) está certa e não muda nada.%
\item%
\textbf{c)} Mesma conclusão e mesmos números: P = 350$\cdot$240/720 = 116,7 kW; Q\_água = (50/60)$\cdot$4000$\cdot$10 = 33,3 kW; $\dot W$\_B = 150 - 45 - 33,3 = 71,7 kW $\Rightarrow$ 215 residências. A resolução diz que o caso com água é compatível com a 2ª lei, o que bate com o gabarito: $\Delta$$\dot S$ = -150/900 + 45/450 + 33,3/300 = +0,044 kW/K. A resolução acrescenta duas coisas corretas: (i) sem a água, a 2ª lei dá $\dot W$\_B $\leq$ 105 - 300(1/6 - 1/10) = 85 kW; (ii) $\dot W$ = 116,7 kW exigiria Q2 = -11,7 kW, com $\dot S$\_ger = -0,106 kW/K. Essas contas confirmam a conclusão e não a contradizem.%
\end{itemize}%
\paragraph{Pontos de dúvida}%
\begin{itemize}%
\item%
\textit{(statement)} Tanto a resolução quanto o gabarito leem o 'Q/2' do esquema da empresa B como metade do calor rejeitado por A, isto é, 45 kW a 450 K. A leitura bate com o texto ('metade do calor associado ao correspondente da empresa A'), mas é uma hipótese de leitura do esquema.%
\item%
\textit{(solution)} O gabarito considera que todo o Q2 vai para a água (Q2 = 33,3 kW); a resolução usa Q2 $\geq$ 33,3 kW. Dá o mesmo limite, $\dot W$\_B $\leq$ 71,7 kW (cerca de 215 residências). A resolução diz que 116,7 kW 'excede o máximo de 85 kW permitido pela 2ª lei', e o gabarito diz que a configuração com água 'não viola a 2ª lei'. Não há contradição, porque as frases tratam de situações diferentes. Ainda assim, o corretor precisa perceber que a resolução também afirma, como o gabarito, que o caso de 71,7 kW com água é compatível com a 2ª lei.%
\item%
\textit{(key)} O gabarito escreve $\Delta$$\dot S$\_universo = +0,04 kW/K; o valor exato é 0,0444 kW/K. É só arredondamento.%
\end{itemize}%
\section*{Questão 6}%
\href{https://poliedroresolve.sistemapoliedro.com.br/vestibulares/ime/2027/ime-2-fase-dia-2-27-10-2026-2-fase-dia-2/dissertativa-6-fisica-2-fase-dia-2-ime-2027}{Resolução do gabarito}\par\smallskip%
\begin{longtable}{@{}p{0.06\textwidth} p{0.33\textwidth} p{0.33\textwidth} p{0.1\textwidth} p{0.08\textwidth}@{}}\toprule Item & Resolução & Gabarito & Veredito & Certo \\ \midrule
a & v = 68 m/s & v = 68 m/s & igual & ambos \\
b & R = 5$\surd$3 m $\approx$ 8,66 m (cos $\alpha$ = 1/2, $\alpha$ = 60$^\circ$) & $\theta$ = 60$^\circ$ (cos $\theta$ = 1/2); R = 5$\surd$3 m & igual & ambos \\
c & T $\approx$ 1,07$\times$10$^3$ N ($\approx$ 1068 N), com $\theta$ = 60$^\circ$ entre o fio e a vertical & $\alpha$ = 60$^\circ$; T $\cong$ 1,07$\cdot$10$^3$ N & igual & ambos \\
d & E $\approx$ 551,2 N; V\_sub $\approx$ 5,5$\times$10$^{-}$$^2$ m$^3$ ($\approx$ 55 L) & E $\approx$ 551 N; V $\cong$ 55$\cdot$10$^{-}$$^3$ m$^3$ & equivalente & ambos \\
\bottomrule\end{longtable}%
\begin{itemize}%
\item%
\textbf{a)} Em A, no sentido anti{-}horário, a velocidade aponta para a direita, ao longo de AO e em direção a O. Então 5/4 = 340/(340 - v), e daí v = 68 m/s. Conferido numericamente.%
\item%
\textbf{b)} Em B a fonte se afasta e só conta a componente v$\cdot$d/$\surd$(d$^2$ + 4R$^2$) da velocidade. Pela fórmula, 340 + v$\cdot$cos = 374, então v$\cdot$cos = 34 e cos = 1/2. A resolução chega a $\surd$(d$^2$ + 4R$^2$) = 2d; o gabarito usa tg 60$^\circ$ = 2R/d. Os dois caminhos são equivalentes e dão R = 5$\surd$3 $\approx$ 8,66 m. Os nomes dos ângulos ($\alpha$ na resolução, $\theta$ no gabarito) estão trocados, mas isso não muda nada.%
\item%
\textbf{c)} Temos sen $\theta$ = R/L = $\surd$3/2, logo $\theta$ = 60$^\circ$. De T$\cdot$sen $\theta$ = m$\cdot$v$^2$/R sai T = $\surd$3$\cdot$68$^2$/(5$\surd$3$\cdot$$\surd$3/2) $\approx$ 1067,9 N. O cálculo com a constante dada (2$\cdot$0,8$\cdot$667,4 $\approx$ 1067,8 N) também confere.%
\item%
\textbf{d)} Na vertical, E = mg + T$\cdot$cos $\theta$ = 17,32 + 533,9 $\approx$ 551,25 N. Então V = E/($\rho$g) $\approx$ 0,0551 m$^3$. O valor é o mesmo; só a notação muda (5,5$\times$10$^{-}$$^2$ m$^3$ e 55$\cdot$10$^{-}$$^3$ m$^3$).%
\end{itemize}%
\paragraph{Pontos de dúvida}%
\begin{itemize}%
\item%
\textit{(statement)} O enunciado diz que o observador está 'na superfície da piscina' e dá apenas a velocidade do som no ar. A resolução e o gabarito supõem que o som chega a O pelo ar, com ar e observador em repouso, e que a fonte está na superfície. A hipótese é razoável e coerente com os dados, mas não está escrita de forma explícita.%
\item%
\textit{(statement)} Algumas informações saem só da figura: o sentido da velocidade em A (para a direita, em direção a O), a tangência de AO à circunferência em A e a fixação do fio no fundo, sobre o eixo vertical. A resolução declara essas leituras e elas conferem com o original. A profundidade implícita, L$\cdot$cos 60$^\circ$ = 5 m, é coerente.%
\item%
\textit{(key)} O gabarito usa $\theta$ para o ângulo de Doppler no item b e $\alpha$ para o ângulo do fio no item c. A resolução faz o contrário. É só convenção de nome e não afeta os resultados.%
\end{itemize}%
\section*{Questão 7}%
\href{https://poliedroresolve.sistemapoliedro.com.br/vestibulares/ime/2027/ime-2-fase-dia-2-27-10-2026-2-fase-dia-2/dissertativa-7-fisica-2-fase-dia-2-ime-2027}{Resolução do gabarito}\par\smallskip%
\begin{longtable}{@{}p{0.06\textwidth} p{0.33\textwidth} p{0.33\textwidth} p{0.1\textwidth} p{0.08\textwidth}@{}}\toprule Item & Resolução & Gabarito & Veredito & Certo \\ \midrule
a & d = 0,45 m & d = 0,45 m & igual & ambos \\
b & E\_el/|$\Delta$E\_pg| = 29/18 $\approx$ 1,6 (kx = 58/9 $\approx$ 6,44 N) & kx $\cong$ 6,44 N; U\_el/$\Delta$U\_g $\cong$ -1,6 & equivalente & ambos \\
c & v\_P2 = $\surd$2 $\approx$ 1,41 m/s, ao atingir o vértice C, sentido de B para C & v2 = $\surd$2 m/s & igual & ambos \\
\bottomrule\end{longtable}%
\begin{itemize}%
\item%
\textbf{a)} Refiz a conta: K\_e q1q2 = 0,9 N$\cdot$m$^2$ e mg sen30$^\circ$ = 2 N. A energia dá 2(1-d) = 0,9(1-d)/d, logo d = 0,45 m (a raiz d = 1 é a posição inicial). Também conferi a coerência física. Em d0 a repulsão vale 0,9 N, menor que 2 N, então P1 desce. O equilíbrio fica em $\surd$0,45 $\approx$ 0,67 m, e o ponto de retorno 0,45 m está abaixo dele, como deve ser.%
\item%
\textbf{b)} A física e os números são os mesmos. Sobre P2 agem, no sentido de descida, o peso tangencial (2 N) e a repulsão de P1 (0,9/0,45$^2$ $\approx$ 4,44 N). A mola, presa em C, equilibra os dois: kx $\approx$ 6,44 N. Temos U\_el = $\frac12$kx$^2$ e $\Delta$U\_g = -mg sen30$^\circ$$\cdot$x = -2x, porque P2 desce x. Logo U\_el/$\Delta$U\_g = -kx/4 = -29/18 $\approx$ -1,61. O gabarito dá a razão literal pedida, com sinal negativo. A resolução usa explicitamente |$\Delta$E\_pg| e chega a +1,61. As duas afirmações são verdadeiras e têm o mesmo módulo. Só o gabarito traz o sinal da variação, que é negativa, e por isso responde de forma mais literal ao enunciado.%
\item%
\textbf{c)} Após o giro, BC fica horizontal e as partículas partem do repouso. Pela conservação do momento, com massas iguais, os deslocamentos são iguais e opostos. P2 percorre 0,15 m até C enquanto P1 percorre 0,15 m e ainda está a 0,40 m de B, então C é o primeiro vértice atingido. A distância passa de 0,45 m para 0,75 m. A energia liberada é 0,9(1/0,45 - 1/0,75) = 0,8 J = 0,4v$^2$, o que dá v = $\surd$2 $\approx$ 1,414 m/s. Resultado confirmado numericamente.%
\end{itemize}%
\paragraph{Pontos de dúvida}%
\begin{itemize}%
\item%
\textit{(solution)} Item b: o enunciado pede a razão entre U\_el e a "variação da energia potencial gravitacional de P2". Essa variação é negativa, porque P2 desce. O gabarito responde $\cong$ -1,6. A resolução responde +1,6 e deixa explícito que usa |$\Delta$E\_pg|. As duas contas são coerentes entre si, mas um corretor rigoroso pode descontar pela falta do sinal. Vale levar ao professor para saber se o sinal é exigido.%
\item%
\textit{(statement)} Item b: a resolução e o gabarito supõem que, com k muito alto, P2 está em equilíbrio quase estático no instante da aproximação máxima, ou seja, kx = peso tangencial + força elétrica. O enunciado não diz isso. Também não diz que P2 começa em equilíbrio: a mola está "inicialmente relaxada", mas P2 já sofre 2 N + 0,9 N. Isso é aceitável porque x $\to$ 0, mas é uma hipótese implícita.%
\item%
\textit{(statement)} O comprimento BC = 1,15 m (d0 + L0) e as posições após o giro (P1 a 0,55 m de B, P2 a 0,15 m de C) vêm da leitura da Figura 2, com P1 no vértice B e P2 a L0 de C, desprezando a compressão da mola. A resolução e o gabarito adotam a mesma leitura.%
\item%
\textit{(statement)} O item c diz "após a condição da Figura 2". Pelo contexto, refere{-}se ao movimento depois do giro e da retirada da mola, isto é, à situação da Figura 3. A redação é ambígua, mas não muda a resposta.%
\end{itemize}%
\section*{Questão 8}%
\href{https://poliedroresolve.sistemapoliedro.com.br/vestibulares/ime/2027/ime-2-fase-dia-2-27-10-2026-2-fase-dia-2/dissertativa-8-fisica-2-fase-dia-2-ime-2027}{Resolução do gabarito}\par\smallskip%
\begin{longtable}{@{}p{0.06\textwidth} p{0.33\textwidth} p{0.33\textwidth} p{0.1\textwidth} p{0.08\textwidth}@{}}\toprule Item & Resolução & Gabarito & Veredito & Certo \\ \midrule
a & Afastando{-}se da rede (o número de máximos visíveis só cai se D aumentar, pois y\_m = D$\cdot$tg$\theta$\_m) & O anteparo está se afastando da rede & igual & ambos \\
b & D(t1) = 0,2$\cdot$$\surd$7 $\approx$ 0,53 m & D1 = 0,53 m & igual & ambos \\
c & v $\approx$ 0,25 m/s (0,254 m/s com os valores da tabela) & v = 0,254 m/s & equivalente & ambos \\
d & $\Delta$s $\approx$ 1,79 m (até D = 0,6$\cdot$$\surd$15 $\approx$ 2,32 m) & $\Delta$s = 1,792 m & equivalente & ambos \\
\bottomrule\end{longtable}%
\begin{itemize}%
\item%
\textbf{a)} Os ângulos $\theta$\_m são fixos (sen $\theta$\_m = m/4). Com y\_m = D$\cdot$tg$\theta$\_m, a queda 7$\to$5$\to$3 exige que D cresça. Resolução e gabarito usam a mesma justificativa.%
\item%
\textbf{b)} Em t1 os máximos $\pm$3 estão na borda: D$\cdot$(3/$\surd$7) = 0,6, logo D = 0,2$\surd$7. Pela tabela dá 0,530 m; o valor exato é 0,529 m. Mesma expressão e mesmo valor nos dois.%
\item%
\textbf{c)} D2 = 0,6$\surd$3 = 1,038 m (pela tabela) e v = (1,038 - 0,530)/2 = 0,254 m/s. Sem a tabela, v = 0,255 m/s. A resolução arredonda para 0,25 m/s, mas cita explicitamente 0,254 m/s. A diferença é só de arredondamento.%
\item%
\textbf{d)} Os máximos $\pm$1 saem em D = 0,6$\surd$15 = 0,6$\cdot$3,87 = 2,322 m, logo $\Delta$s = 2,322 - 0,530 = 1,792 m (o valor exato é 1,795 m). A resolução dá 1,79 m, a mesma resposta arredondada.%
\end{itemize}%
\paragraph{Pontos de dúvida}%
\begin{itemize}%
\item%
\textit{(statement)} A contagem de 7 máximos supõe que a ordem m = $\pm$4 (sen $\theta$ = 1, $\theta$ = 90$^\circ$) é rasante e nunca chega ao anteparo. Também trata a rede como pontual e as franjas como pontos, de modo que o máximo é visto enquanto y $\leq$ L/2. São hipóteses razoáveis e coerentes com o '7' do enunciado. Resolução e gabarito adotam as duas.%
\item%
\textit{(solution)} Nos itens c e d, os valores com a tabela de raízes (0,254 m/s e 1,792 m) diferem um pouco dos exatos (0,255 m/s e 1,795 m). A resolução arredonda para 0,25 m/s e 1,79 m. Essa diferença de arredondamento não deve afetar a nota.%
\end{itemize}%
\section*{Questão 9}%
\href{https://poliedroresolve.sistemapoliedro.com.br/vestibulares/ime/2027/ime-2-fase-dia-2-27-10-2026-2-fase-dia-2/dissertativa-9-fisica-2-fase-dia-2-ime-2027}{Resolução do gabarito}\par\smallskip%
\begin{longtable}{@{}p{0.06\textwidth} p{0.33\textwidth} p{0.33\textwidth} p{0.1\textwidth} p{0.08\textwidth}@{}}\toprule Item & Resolução & Gabarito & Veredito & Certo \\ \midrule
a & u = c($\surd$3h - L)/($\surd$3L - h), com t = ($\surd$3L - h)/c e R = $\surd$3h - L no impacto & u = c($\surd$3$\cdot$h - L)/($\surd$3$\cdot$L - h) & igual & ambos \\
b & h < L < $\surd$3h (ou seja, 1 < L/h < $\surd$3) & L/$\surd$3 < h < L & equivalente & ambos \\
\bottomrule\end{longtable}%
\begin{itemize}%
\item%
\textbf{a)} Mesma expressão. Refiz a conta: reflexão de d=($\surd$3/2, 1/2) para d'=(0,-1) dá a normal externa n=(-1/2, -$\surd$3/2) (incidência de 30$^\circ$). Igualando P = C + ut$\cdot$n a P = ct$\cdot$d, saem $\surd$3ct/2 + ut/2 = L e ct/2 + $\surd$3ut/2 = h, logo ut = $\surd$3h - L e ct = $\surd$3L - h. Conferi em Python para (h,L) = (1;1,5), (1;1,2) e (1;1,7): calculando a menor raiz de |ct$\cdot$d - C| = ut (primeiro contato), o raio refletido sai exatamente (0, -1).%
\item%
\textbf{b)} São a mesma condição escrita de outro jeito: h < L < $\surd$3h $\Leftrightarrow$ (h < L e h > L/$\surd$3) $\Leftrightarrow$ L/$\surd$3 < h < L. As duas resoluções impõem ut > 0 (L < $\surd$3h), ct > 0 (h < $\surd$3L, que é automática quando h < L) e u < c ($\surd$3h - L < $\surd$3L - h $\Leftrightarrow$ h < L). A resolução ainda mostra que a raiz encontrada é o primeiro contato, o que o gabarito não faz.%
\end{itemize}%
\paragraph{Pontos de dúvida}%
\begin{itemize}%
\item%
\textit{(statement)} O enunciado não diz que a esfera precisa crescer de fato (u > 0). Tanto a resolução quanto o gabarito usam desigualdade estrita, excluindo u = 0 (L = $\surd$3h), caso em que a esfera vira um ponto sobre a reta e não há reflexão. Faz sentido físico, mas é uma hipótese que nenhum dos dois declara.%
\item%
\textit{(key)} O gabarito não verifica que a solução corresponde ao primeiro contato do fóton com a esfera (menor raiz de |ct$\cdot$d - C| = ut) nem que o fóton chega por fora. A resolução faz essa verificação (f' < 0 no impacto). Isso não muda a resposta.%
\end{itemize}%
\section*{Questão 10}%
\href{https://poliedroresolve.sistemapoliedro.com.br/vestibulares/ime/2027/ime-2-fase-dia-2-27-10-2026-2-fase-dia-2/dissertativa-10-fisica-2-fase-dia-2-ime-2027}{Resolução do gabarito}\par\smallskip%
\begin{longtable}{@{}p{0.06\textwidth} p{0.33\textwidth} p{0.33\textwidth} p{0.1\textwidth} p{0.08\textwidth}@{}}\toprule Item & Resolução & Gabarito & Veredito & Certo \\ \midrule
a & Somente a situação 2. Nas situações 1 e 4 (correntes opostas), B\_z $\neq$ 0 no eixo e F\_x tem sempre o mesmo sinal, logo não há restauração. Na situação 3, F\_x = +qvB'x afasta a partícula do eixo (instável). Na situação 2, F\_x = -qvB'x é restauradora. & Somente a situação 2 (1 e 4: B $\neq$ 0 no eixo; 3: a força afasta; 2: a força restaura). & igual & ambos \\
b & p = $\pi$v/$\omega$ $\approx$ ($\pi$d/2)$\cdot$$\surd$($\pi$ m v/($\mu$0 i q)), com $\omega$$^2$ = 4$\mu$0 i q v/($\pi$ m d$^2$) (T = 2$\pi$/$\omega$) & T = $\pi$d$\cdot$$\surd$(m$\pi$/($\mu$0 i q v)); p = ($\pi$d/2)$\cdot$$\surd$($\pi$ m v/($\mu$0 i q)) & igual & ambos \\
\bottomrule\end{longtable}%
\begin{itemize}%
\item%
\textbf{a)} Conferi com a figura original da p. 10. Com os fios em x = $\pm$d/2 e as correntes para cima, B\_z = ($\mu$0 i/2$\pi$){[}1/(d/2-x) - 1/(d/2+x){]} $\approx$ 4$\mu$0 i x/($\pi$d$^2$), e F\_x = q v\_y B\_z. Na situação 2 (v\_y = -v), F\_x = -(4$\mu$0 i q v/($\pi$d$^2$))x, que é restauradora. Na situação 3 a força tem o sinal de x e afasta a partícula. Com correntes opostas (1 e 4), os campos se somam entre os fios e B\_z $\approx$ -B0(1+4x$^2$/d$^2$) nunca se anula. Assim, a força lateral já existe no eixo e tem sentido único. A resolução também acerta para qual fio a partícula é desviada (o de corrente paralela ao movimento). As justificativas coincidem com as do gabarito.%
\item%
\textbf{b)} k = 4$\mu$0 i q v/($\pi$d$^2$) e T = 2$\pi$$\surd$(m/k) = $\pi$d$\surd$($\pi$m/($\mu$0 i q v)), o mesmo período do gabarito. Daí p = vT/2 = ($\pi$d/2)$\surd$($\pi$mv/($\mu$0 i q)). Conferi com python: com os dados do teste da resolução, $\pi$v/$\omega$ e vT/2 dão os dois 0,80261 m. Uma integração RK4 com o campo exato dos dois fios (x0 = 2 mm) deu 0,80212 m entre cruzamentos consecutivos, o mesmo valor que a resolução informa. O raio R = mv/(qB0) = 1,305 m das situações 1 e 4 também confere. A resolução não escreve o período T como expressão fechada, só como T = 2$\pi$/$\omega$ com $\omega$ dado explicitamente. Isso é equivalente, e o que se pede é p.%
\end{itemize}%
\paragraph{Pontos de dúvida}%
\begin{itemize}%
\item%
\textit{(statement)} O enunciado define p como a distância entre dois cruzamentos do eixo y, e a figura mostra cruzamentos consecutivos, o que dá p = vT/2. Quem lesse p como o comprimento de onda completo (vT) chegaria a um valor 2 vezes maior. A resolução e o gabarito adotam a mesma leitura, a de meio período, que é a mais coerente com a figura.%
\item%
\textit{(solution)} A resolução despreza a gravidade, que não é dada, e admite fios eletricamente neutros. São hipóteses implícitas razoáveis, que ela declara, e o gabarito usa as mesmas.%
\end{itemize}%
\end{document}